\subsection{The vacuity gap and self-verification}
\label{sec:vacuity}

The kernel is a strong guarantee, but it is important to be precise about what it
guarantees. The kernel certifies that a proof term proves \emph{its stated
theorem}. It does not certify that the stated theorem is the one you meant, nor
that the theorem says anything at all. A theorem can be true and yet
\emph{vacuous}---$X = X$ proves nothing, and a more subtle relative can hide the
same emptiness behind real-looking algebra. The kernel will happily accept it.
This is the one gap Telperion is built around, and we would rather name it
loudly than paper over it.

\subsection{Making emitters declare themselves}

Every emitter must declare a \emph{sensitivity stance}, and a test refuses to let
one ship without it:

\begin{itemize}
  \item \textbf{Certificate-sensitive.} The emitted theorem carries a corruptible
    witness---a specific numeric or algebraic identity that is doing the work. For
    these, a wrong witness must break the proof. We check exactly that with a
    \emph{negative control}: forge the witness and confirm the kernel
    \emph{rejects} the forgery. If corrupting the certificate still compiles, the
    certificate was not load-bearing, and the emitter is refused.
  \item \textbf{Structurally non-vacuous.} The theorem is a positivity,
    decidability, finite-cover, or hypothesis-gated statement with no separable
    corruptible witness. Here a structural non-vacuity check---run at emit time,
    refusing reflexive and trivial statements---is the right guard, and the
    contributor writes one line saying why.
\end{itemize}

The distinction is not bureaucracy; it is the difference between ``the kernel
accepted a proof'' and ``the kernel accepted a proof that could have failed.''
Only the second is evidence.

\subsection{A second, independent kernel}

For an even stronger check, Telperion integrates an independent judge: the
Comparator \cite{comparator}, drawn from OpenAI's ten-proofs project
\cite{tenproofs}, together with \texttt{nanoda}, a from-scratch reimplementation
of the Lean type checker in Rust \cite{nanoda}. The Comparator confirms that an
emitted proof proves \emph{exactly} the challenge statement using only a
whitelisted set of axioms. Two independently written checkers agreeing is a
meaningfully higher bar than one.

\subsection{The gap, encountered in the wild}

This is not a hypothetical hazard for us; it visited this very campaign. Midway
through the counting-formula program (\cref{sec:rh}), a planning claim---``this
entire function carries the nontrivial zeros''---was simply false, while every
kernel theorem built around it remained true: the reflection identities were
correct statements about the wrong object. The kernel could not have caught it,
because the false claim lived in docstrings and plans, not in any theorem. It
was caught one step later by checking the premise against Mathlib's own
defining identity, the detour was relabeled rather than hidden, and the durable
fix is the one this section recommends in general: the successor object's zero
set was not asserted but \emph{proved}, as an iff, in the axiom guard. The
kernel guards statements; the discipline must guard what the statements are
about.

\subsection{The gap, encountered a second time---and caught by the machinery}

A harder instance arrived during the Turing-ladder campaign (\cref{sec:rh}),
and this time the discipline of this section, not luck, caught it. The original
winding-route band certificates carried an edge-hypothesis bundle of the shape
``for every decomposition of the boundary data satisfying $P$, the interior
zeros are confined''---universally quantified over a decomposition that the
intended Arb data was supposed to instantiate. Auditing the shape before
scaling the campaign, we asked the sensitivity question this section mandates:
\emph{what else satisfies the premise?} The answer was: junk. A decomposition
padding the zero multiset with a phantom point of multiplicity zero, or
absorbing an off-ball point into the error term, satisfies $P$ while violating
the confinement clause---so the bundled hypothesis was false for every band,
the per-band theorems true but \emph{vacuous} as certificates, and the
milestone claims built on them empty. The kernel had accepted every one of
them, because every one of them was true. Three consequences followed, each the
recommendation of this section made concrete. First, the affected milestone
claims were \emph{suspended publicly}, in the pull request and the campaign
document, the day the vacuity was found. Second, the replacement certificate
shape was made \emph{canonical and audited}: every re-based band instantiates
one fixed Prop (\texttt{TuringBand.BandStatement}) whose shape is pinned by
hand-audited kernel code, so an emitter can vary only parameters that are
numerically cross-checked, and a kernel statement-match gate plus
forge-the-witness negative controls run over every band in continuous
verification. Third, the entire ladder---every zero of every prior
milestone---was re-verified on the new shape before any claim was reinstated;
the re-based census agrees with the rigorous count $N(T)$ exactly at every
thousand-block. The episode cost one day. Scaling a quarter-million-zero
verification on a vacuous certificate shape would have cost the paper.

\subsection{The floor we cannot get under}

Honesty requires stating the limit, too. No amount of machinery decides whether a
\emph{formal} statement faithfully captures the \emph{informal} claim a human had
in mind; that correspondence lives outside any formal system, and by the usual
self-reference obstructions (L\"ob, G\"odel) it cannot be internalized. So the
irreducible trusted component is not Telperion and not even the kernel---it is the
human reading of the theorem statements. Everything in this section shrinks the
untrusted surface down to that floor; nothing can remove the floor itself. We
take that as a reason to keep the statements short, legible, and few, not as a
reason to pretend the floor is not there.
